% MCSPMAGCONE1 2026-10-17 21:45 Asia/Seoul
\documentclass[11pt]{article}
\usepackage{amsmath,amssymb,amsthm,hyperref}
\newtheorem{theorem}{Theorem}
\newtheorem{lemma}[theorem]{Lemma}
\newtheorem{corollary}[theorem]{Corollary}
\title{MCSPMAGCONE1: Applying the Project's Asymptotic-Cone and Locality Machinery to Hardness Magnification}
\author{P versus NP project}
\date{17 October 2026}
\begin{document}\maketitle
\noindent\textbf{Status.} The sequence-magnification corollary and the local-shattering lemma are PROVED/EXACT from the cited magnification theorems and elementary circuit interpolation. The proposed streaming/version-space and HD routes are CANDIDATE. No MCSP lower bound crossing the magnification threshold is proved; $P\ne NP$ remains OPEN.

\section{IDEAS PROPOSED BEFORE LITERATURE}
A: replace one exact magnification parameter by a sequence $\beta_j\downarrow0$ and a fixed lower-bound excess. B: apply the project's local-moment/HD machinery to the low-circuit truth-table family. C: apply erasure/continuation-DAG ideas to search-MCSP streaming. D: use the formula/branching-program magnification frontiers only as secondary targets, not as the final P-vs-NP implication unless the conclusion is $NP\not\subseteq P/poly$.

\section{LITERATURE CHECK AND REVISIONS}
Oliveira--Pich--Santhanam prove that there is a universal $c\ge1$ such that if some $\varepsilon>0$ satisfies, for every sufficiently small $\beta>0$,
\[\mathrm{GapMCSP}[2^{\beta n}/(cn),2^{\beta n}]\notin\mathrm{Circuit}[N^{1+\varepsilon}],\quad N=2^n,\]
then $NP\not\subseteq P/poly$. Their proof under the contrary assumption gives almost-linear circuits for every sufficiently small $\beta$, via anti-checkers.

McKay--Murray--Williams prove that if search-MCSP with parameter $s(n)$ lacks a one-pass streaming algorithm using $\operatorname{poly}(s)$ space and $\operatorname{poly}(s)$ update time, then $P\ne NP$. Cheraghchi--Hirahara--Myrisiotis--Yoshida restate a corollary: for some constant $\mu>0$, an $N^{1.01}$ one-tape lower bound for $\mathrm{MCSP}[2^{\mu n}]$ suffices for $P\ne NP$, while their unconditional $N^{1.99}$ lower bound holds at a larger parameter $\mu_2>\mu_1$.

Chen--Hirahara--Oliveira--Pich--Rajgopal--Santhanam identify the locality barrier: magnification targets already have very efficient circuits with small-fan-in oracle gates, so lower-bound methods stable under such local oracles are unlikely to cross the magnification frontier.

\section{CONSTRUCTION OR PROOF ATTEMPT}
\begin{theorem}[Sequence magnification for Gap-MCSP]
Let $c$ be the universal constant of the Oliveira--Pich--Santhanam theorem. If there exist $\varepsilon>0$ and a sequence $\beta_j\downarrow0$ such that for every $j$
\[\mathrm{GapMCSP}[2^{\beta_j n}/(cn),2^{\beta_j n}]\notin\mathrm{Circuit}[N^{1+\varepsilon}],\]
then $NP\not\subseteq P/poly$, hence $P\ne NP$.
\end{theorem}
\begin{proof}
Assume $NP\subseteq P/poly$. The anti-checker proof of the magnification theorem yields $N^{1+\varepsilon}$ circuits for every sufficiently small $\beta>0$. Choose $j$ with $\beta_j$ below that threshold, contradicting the assumed lower bound.
\end{proof}

This is the project's limit-space formulation: we only need a fixed vertical gap $\varepsilon>0$ along any parameter sequence approaching the origin. No exact optimal $\beta$ is required.

\begin{corollary}[Streaming cone criterion]
Suppose there is a constant $\delta>0$ and a sequence $\beta_j\downarrow0$ such that search-MCSP at $s=N^{\beta_j}$ requires either space or update time greater than $N^\delta$ for every one-pass streaming algorithm. Then $P\ne NP$.
\end{corollary}
\begin{proof}
Under $P=NP$, McKay--Murray--Williams give one fixed polynomial $p$ such that search-MCSP has space and update time $p(s)$. Write $p(s)\le s^d$ for fixed $d$. For $j$ with $d\beta_j<\delta$, $p(N^{\beta_j})<N^\delta$, contradiction.
\end{proof}

An analogous one-tape cone has upper exponent $1+O(\beta)$ under $P=NP$; therefore a fixed lower exponent $1+\delta$ along $\beta_j\downarrow0$ suffices. Existing work already supplies exponent $1.99$ at a larger fixed parameter; the missing progress is chiefly moving the parameter toward zero, not increasing the time exponent.

\begin{lemma}[Local shattering barrier for MCSP]
Fix $r$ distinct inputs $x_1,\ldots,x_r\in\{0,1\}^n$ and arbitrary labels $b_1,\ldots,b_r$. There is a Boolean circuit of size $O(rn)$ agreeing with all these labels.
\end{lemma}
\begin{proof}
For every $i$ with $b_i=1$, build the equality indicator $[x=x_i]$ using $O(n)$ gates and OR these indicators; if all labels are zero use the constant-zero circuit.
\end{proof}

Therefore the class of truth tables computable by circuits of size $s$ shatters every coordinate set of size $r\le\Theta(s/n)$. For $s=N^\beta$, any observation of fewer than about $N^\beta/\log N$ truth-table coordinates can realize every local bit pattern inside the MCSP YES class. Thus the project's clause-moment, sunflower, and low-order HD methods do not transfer directly at logarithmic order. A successful MCSP application must reach observation order at least $N^\beta/\operatorname{polylog}N$, or exploit nonlocal minimization structure.

This is a concrete instance of the published locality barrier and retires a naive transfer of the exact-SAT slice moment method.

\paragraph{Streaming/version-space reformulation.}
After reading truth-table positions $(x_i,b_i)$, define
\[\mathcal V_i=\{C:|C|\le s,\ C(x_j)=b_j\text{ for all }j\le i\}.\]
A search-MCSP streaming algorithm is an online compressed representation of this monotonically shrinking version space that must finally output a member if $\mathcal V_N\ne\varnothing$. The project's signed-DAG/continuation-equivalence machinery is relevant here: histories may merge exactly when they induce equivalent future version spaces. Pure state counting is insufficient because one circuit description already uses only $O(s\log s)$ bits, within $\operatorname{poly}(s)$ space. The target must be a computational lower bound on maintaining or updating the quotient, not an information-volume lower bound.

\paragraph{Lipschitz/Hamming control.}
If truth tables $f,g$ differ on $h$ entries, then
\[|CC(f)-CC(g)|\le O(nh),\]
because one may correct each differing truth-table position with a point-indicator circuit of size $O(n)$. Hence the Gap-MCSP YES and NO regions with circuit-size gap $s_2-s_1$ are separated in Hamming distance by at least $\Omega((s_2-s_1)/n)$. This is useful geometry but does not by itself lower-bound decision circuits: a Hamming-separated predicate can still be easy. It should be treated as a structural constraint, not a lower bound.

\paragraph{Formula frontier.}
Known local-PRG lower bounds already give $N^{3-o(1)}$ De Morgan formula lower bounds and $N^{2-o(1)}$ arbitrary-formula/branching-program lower bounds for MCSP. Magnification theorems ask for a fixed $\varepsilon$ improvement past such frontier exponents in several variants. This is where the project's HD/KW/infinite-series methods may be most technically compatible. However, conclusions such as $NP\not\subseteq NC^1$ or $\oplus P\not\subseteq NC^1$ do not by themselves imply $P\ne NP$; the final spine should retain the $NP\not\subseteq P/poly$ or direct streaming-to-$P\ne NP$ magnification endpoint.

\section{TESTS AND OUTCOME}
The best-fit proof spine is now two-pronged. Route I: prove a fixed $N^{1+\varepsilon}$ circuit lower bound for Gap-MCSP along any $\beta_j\to0$; the sequence magnification theorem finishes $P\ne NP$. Route II: prove a fixed $N^\delta$ one-pass streaming space/update lower bound for search-MCSP along any $\beta_j\to0$; McKay--Murray--Williams finishes $P\ne NP$. The second route matches the project's erasure, accounting, and continuation-DAG work more closely, while the first has the cleanest nonuniform conclusion.

The principal barrier is not the exponent itself but uniformity as the compression parameter tends to zero. Existing MCSP lower bounds are strong in exponent but live at larger parameter regimes. The new target is therefore an asymptotic lower-bound region with a positive intercept at $\beta=0$.

Root main.tex was not modified.
\end{document}
% FINAL_VERIFICATION_SENTINEL::MAIN::MCSPMAGCONE1::2026-10-17_2145::P_VS_NP_OPEN